\documentclass[12pt]{article}

\usepackage{amsmath}
\usepackage{amssymb}
\usepackage{graphicx}
\graphicspath{{figures/}{figures/seed_selection/}}
\usepackage{hyperref}
\usepackage[margin=1in]{geometry}
\usepackage{xcolor}
\usepackage{booktabs}
\usepackage{tcolorbox}
\tcbuselibrary{breakable,skins}
\usepackage{natbib}
\bibliographystyle{aasjournal}

\newcommand{\bk}{\mathbf{k}}
\newcommand{\bx}{\mathbf{x}}
\newcommand{\br}{\mathbf{r}}
\newcommand{\avg}[1]{\left\langle #1 \right\rangle}
\newcommand{\Var}{\mathrm{Var}}
\newcommand{\Cov}{\mathrm{Cov}}
\newcommand{\Mpch}{\ \mathrm{Mpc}/h}

\title{Selecting a Random Seed on the Correlation Function\\
       of a Subvolume}
\author{}
\date{\today}

\begin{document}

\maketitle

\begin{tcolorbox}[colback=yellow!10!white, colframe=orange!75!black,
                  title={\bfseries Disclaimer}]
These notes were compiled by an AI assistant (Claude, Anthropic) and have
\textbf{not been reviewed by a human domain expert}.  Cross-check key
equations and conclusions before relying on them.
\end{tcolorbox}

\begin{abstract}
\noindent
The correlation function $\xi(r)$ measured in a $125 \times 125 \times
1000\Mpch$ subvolume of a $1\ \mathrm{Gpc}/h$ box departs from linear theory at
separations beyond a few tens of $\mathrm{Mpc}/h$.  One proposed remedy is to
search over random seeds and adopt the one whose subvolume $\xi$ follows theory
most closely.  We show that this replaces the intended ensemble by the
conditional ensemble given the search criterion, and we quantify the
consequences with 600 Gaussian realizations of the same volume and shape.
Retaining the best fifth of seeds reduces the realization-to-realization
scatter of $\xi$ to $0.38$--$0.64$ of its correct value within the matched
range, and to $0.67$--$0.94$ at separations excluded from the match; ensemble
means at excluded separations move by up to $2.8$ times their own standard
error.  The reduction has no threshold: rejecting only the worst fifth of seeds
already removes $22\%$ of the scatter.  Two further results limit the
usefulness of the procedure.  First, the target is displaced: the integral
constraint places the expectation of the subvolume estimator $0.22\sigma$ below
linear theory at $r = 92\Mpch$, so a seed matching raw theory is atypical by
construction.  Second, the objective is not attainable in practice: the closest
of 600 seeds still departs from theory by $0.19\sigma$ per separation, while an
estimator biased low by $10\%$ would displace the measurement by only
$0.068\sigma$, so the search neither reproduces theory closely nor
distinguishes a correct pipeline from a faulty one.  We recommend comparing
against the expectation of the estimator rather than against raw theory,
reporting the sample-variance band, and, where a variance-reduced realization
is genuinely required, using fixed-amplitude or constrained initial conditions.
Where a seed must be rejected, rejection on the mean density of the subvolume
alone accepts $96\%$ of seeds and removes none of the scatter of $\xi$.
\end{abstract}

\tableofcontents

\section{Introduction}\label{sec:intro}

A cosmological box of side $L = 1000\Mpch$ contains a subvolume of
$125 \times 125 \times 1000\Mpch$, one sixty-fourth of the total volume, which
is to be resimulated at higher resolution.  The two-point correlation function
$\xi(r)$ measured within the subvolume departs from the linear-theory
prediction at separations beyond about $30\Mpch$, while the same measurement
over the full box follows it.

The proposal under examination is to generate initial conditions from many
random seeds and to adopt the seed whose subvolume $\xi(r)$ reproduces the
theoretical curve most closely.  The motivation is that a simulation should
represent the cosmology whose parameters it was given, and that a subvolume
whose clustering disagrees with that cosmology is a poor starting point for a
resimulation.

Section~\ref{sec:conditioning} establishes what such a selection does, first
through an example that can be evaluated exactly and then through the
conditional moments of a Gaussian pair.  Section~\ref{sec:application} applies
the result to the correlation function of a subvolume and identifies a
systematic displacement of the target.  Section~\ref{sec:experiment} quantifies
the effect with 600 realizations.  Sections~\ref{sec:tolerance}
and~\ref{sec:circular} address the dependence on the tolerance and the
circularity of validating a pipeline on a selected seed.
Section~\ref{sec:recommendations} states the alternatives.

Throughout, a \emph{realization} denotes one seed and the field generated from
it; $\avg{\cdot}$ denotes an average over realizations; the \emph{scatter}
$\sigma_Q$ of a quantity $Q$ is the root-mean-square deviation of $Q$ over
realizations; and \emph{theory} denotes the prediction for $\avg{\cdot}$, not
for any individual realization.

\section{Selection as conditioning}\label{sec:conditioning}

\subsection{An exact example}\label{sec:dice}

Let $X_1$ and $X_2$ be the outcomes of two fair dice and $S = X_1 + X_2$.  For
a single die,
\begin{equation}
  \avg{X} = \frac{1}{6}\sum_{i=1}^{6} i = \frac{7}{2}, \qquad
  \Var[X]  = \frac{1}{6}\sum_{i=1}^{6}\Bigl(i - \frac{7}{2}\Bigr)^{2}
           = \frac{35}{12},
  \label{eq:die}
\end{equation}
so $\sigma_X = \sqrt{35/12} = 1.708$.  The dice are independent, hence
$\avg{S} = 7$ and $\Var[S] = 35/6$.

Retain only the rolls with $S = 7$, the value of $\avg{S}$.  Six of the
thirty-six equally likely outcomes satisfy this, and each admits exactly one
value of $X_1$, so
\begin{equation}
  P(X_1 = i \mid S = 7)
    = \frac{P(X_1 = i,\, X_2 = 7-i)}{P(S = 7)}
    = \frac{1/36}{6/36} = \frac{1}{6}.
\end{equation}
The marginal distribution of $X_1$ is therefore unaltered: its mean remains
$7/2$ and its scatter remains $1.708$.  Examined alone, the retained rolls are
indistinguishable from unselected ones.

The joint distribution is altered, and this is where the whole effect
originates, so we set it out in full.  The covariance of two quantities is the
average of the product of their deviations from their own means,
\begin{equation}
  \Cov[A, B] = \avg{(A - \avg{A})(B - \avg{B})},
  \label{eq:cov}
\end{equation}
and the correlation coefficient is that quantity rescaled to a fixed range,
\begin{equation}
  \rho = \frac{\Cov[A, B]}{\sigma_A\,\sigma_B}, \qquad -1 \leq \rho \leq 1 .
  \label{eq:rho}
\end{equation}
A value $\rho = 0$ indicates that one quantity carries no information about the
other; $|\rho| = 1$ indicates that one is an exact linear function of the other.

Before selection the dice are independent: whatever $X_1$ shows, $X_2$ is
equally likely to fall above or below $7/2$, so the product in
equation~\eqref{eq:cov} is positive as often as negative and averages to zero.
After selection $X_2 = 7 - X_1$, so subtracting $7/2$ from both sides gives
\begin{equation}
  X_2 - \frac{7}{2} = -\Bigl(X_1 - \frac{7}{2}\Bigr),
\end{equation}
that is, the second deviation is the negative of the first.  Every product is
then $-(X_1 - 7/2)^2$, negative for every retained roll:
\begin{center}
\begin{tabular}{lccccccc}
  \toprule
  roll & $(1,6)$ & $(2,5)$ & $(3,4)$ & $(4,3)$ & $(5,2)$ & $(6,1)$ & average\\
  \midrule
  $X_1 - 7/2$ & $-2.5$ & $-1.5$ & $-0.5$ & $+0.5$ & $+1.5$ & $+2.5$ & \\
  $X_2 - 7/2$ & $+2.5$ & $+1.5$ & $+0.5$ & $-0.5$ & $-1.5$ & $-2.5$ & \\
  product     & $-6.25$ & $-2.25$ & $-0.25$ & $-0.25$ & $-2.25$ & $-6.25$
              & $-35/12$\\
  \bottomrule
\end{tabular}
\end{center}
The average of the products is $-17.5/6 = -35/12$, which is $-\Var[X_1]$ by
equation~\eqref{eq:die}, and equation~\eqref{eq:rho} then gives $\rho = -1$.
The selection has moved the pair from independence to exact opposition, the
furthest it can go: every unit by which one die runs high, the other runs
equally low.  Any quantity depending on both dice inherits the constraint.  For $Y = X_1 + 2X_2$, independence gives
\begin{equation}
  \avg{Y} = 3\avg{X} = \frac{21}{2}, \qquad
  \Var[Y] = \Var[X_1] + 4\Var[X_2] = \frac{175}{12},
\end{equation}
so $\sigma_Y = 3.819$, whereas substituting $X_2 = 7 - X_1$ gives $Y = 14 -
X_1$ and
\begin{equation}
  \avg{Y \mid S = 7} = \frac{21}{2}, \qquad
  \Var[Y \mid S = 7] = \Var[X_1] = \frac{35}{12},
\end{equation}
so $\sigma_Y$ falls by a factor $\sqrt{5}$, from $3.819$ to $1.708$, at
unchanged mean.  The rule defining the selection was stated in terms of $S$
alone and made no reference to $Y$; the scatter of $Y$ nevertheless fell by more
than half.  The reason is that $Y$ and $S$ are not independent.  Both are built
from the same two dice, and equation~\eqref{eq:rho} gives
$\rho(Y, S) = 3/\sqrt{10} = 0.949$, so constraining $S$ constrains $Y$ with it.
A quantity uncorrelated with $S$ would pass through the selection unaffected.

Selecting on a value other than $\avg{S}$ moves the individual dice as well, and
not only their scatter.  Retaining $S = 7$ left the average of $X_1$ at $7/2$,
because $7$ is the average of $S$.  Retaining the rolls with $S \geq 10$ aims at
a value above the average of $S$, and leaves $(4,6)$, $(5,5)$, $(6,4)$, $(5,6)$,
$(6,5)$, $(6,6)$, for which
\begin{equation}
  \avg{X_1 \mid S \geq 10} = \frac{16}{3} = 5.333, \qquad
  \Var[X_1 \mid S \geq 10] = 29 - \Bigl(\frac{16}{3}\Bigr)^{2}
                           = \frac{5}{9},
\end{equation}
using $\Var[Q] = \avg{Q^2} - \avg{Q}^2$.  The mean of $X_1$ has moved by
$11/6 = 1.07\,\sigma_X$ and its scatter has fallen to $0.44\,\sigma_X$, for a
die that was never examined individually.  Figure~\ref{fig:dice} displays both
selections.

\begin{figure}[htbp]
  \centering
  \includegraphics[width=\textwidth]{seed_dice.pdf}
  \caption{Selection applied to two fair dice.  \emph{(a)} The thirty-six
    equally likely outcomes, with the six satisfying $S = 7$ marked.
    \emph{(b)} Distribution of $Y = X_1 + 2X_2$ before and after retaining
    $S = 7$: the mean is unchanged and the scatter falls by a factor $\sqrt5$,
    although the rule defining the selection made no reference to $Y$.  The two
    are correlated, $\rho(Y,S) = 0.949$, which is why constraining one
    constrains the other.  \emph{(c)} Distribution of
    $X_1$ before and after retaining $S \geq 10$: selecting on a value other
    than $\avg{S}$ displaces the mean of $X_1$ by $1.07\,\sigma_X$ and reduces
    its scatter to $0.44\,\sigma_X$.  Produced by
    \texttt{figures/seed\_selection/plot\_seed\_dice.py}.}
  \label{fig:dice}
\end{figure}

\subsection{Conditional moments of a Gaussian pair}\label{sec:gaussian}

The dice could be counted outcome by outcome.  A density field cannot, so we
carry out the same computation once for a pair of Gaussian quantities, which is
a good approximation for the large scales at issue here.

Let $T$ and $S$ be jointly Gaussian with correlation coefficient $\rho$ as
defined in equation~\eqref{eq:rho}.  It is convenient to measure each in units
of its own scatter, relative to its own mean:
\begin{equation}
  t = \frac{T - \avg{T}}{\sigma_T}, \qquad
  u = \frac{S - \avg{S}}{\sigma_S},
  \label{eq:standardize}
\end{equation}
so that $t$ and $u$ have zero mean and unit scatter, and $\rho = \avg{tu}$.
Their joint probability density is then the standard bivariate normal,
\begin{equation}
  p(t, u) = \frac{1}{2\pi\sqrt{1-\rho^{2}}}
            \exp\left[-\,\frac{t^{2} - 2\rho\,t u + u^{2}}
                              {2\,(1-\rho^{2})}\right].
  \label{eq:bivariate}
\end{equation}
Selecting on $T$ means fixing $t$ and asking what is left of the distribution of
$u$.  That is the conditional density $p(u \mid t) = p(t,u)/p(t)$, where
$p(t) = (2\pi)^{-1/2} e^{-t^{2}/2}$ is the density of $t$ alone.

The quadratic form in equation~\eqref{eq:bivariate} can be rearranged so that
$u$ appears in one place only:
\begin{equation}
  t^{2} - 2\rho\, t u + u^{2}
    = (u - \rho t)^{2} + t^{2}\,(1 - \rho^{2}),
  \label{eq:complete-square}
\end{equation}
which is verified by expanding $(u - \rho t)^{2} = u^{2} - 2\rho t u +
\rho^{2}t^{2}$ and adding $t^{2} - \rho^{2}t^{2}$.  Substituting
equation~\eqref{eq:complete-square} into equation~\eqref{eq:bivariate} splits
the exponential into a factor containing $u$ and a factor containing $t$ alone,
\begin{equation}
  p(t, u) = \underbrace{\frac{1}{\sqrt{2\pi(1-\rho^{2})}}
            \exp\left[-\,\frac{(u - \rho t)^{2}}{2(1-\rho^{2})}\right]}
            _{\textstyle p(u \mid t)}
            \;\times\;
            \underbrace{\frac{1}{\sqrt{2\pi}}\,e^{-t^{2}/2}}_{\textstyle p(t)} ,
\end{equation}
and dividing by $p(t)$ leaves the first factor.  It is a Gaussian in $u$ with
mean $\rho t$ and variance $1 - \rho^{2}$.

Two features of that result are the substance of this note.  The mean of $u$ is
displaced from zero whenever $t \neq 0$, that is whenever the selected value
differs from $\avg{T}$; and the variance is reduced from $1$ to $1-\rho^{2}$
whatever value is selected, including $t = 0$.  Restoring the original units
through equation~\eqref{eq:standardize},
\begin{align}
  \avg{S \mid T = t_0} &= \avg{S}
      + \rho\, \frac{\sigma_S}{\sigma_T}\,\bigl(t_0 - \avg{T}\bigr),
      \label{eq:cond-mean}\\[2pt]
  \Var\bigl[S \mid T = t_0\bigr] &= \sigma_S^{2}\,(1 - \rho^{2}).
      \label{eq:cond-var}
\end{align}
Equation~\eqref{eq:cond-var} states that the scatter of any $S$ correlated with
$T$ is reduced by $\sqrt{1-\rho^2}$ regardless of the value selected;
equation~\eqref{eq:cond-mean} states that a target displaced from $\avg{T}$
displaces $\avg{S}$ in proportion to $\rho$.  These are the two effects
exhibited by the dice.

The dice do not satisfy equations~\eqref{eq:cond-mean}
and~\eqref{eq:cond-var} pointwise, and the discrepancy is instructive.  For $Y$
and $S$ above, $\rho = 3/\sqrt{10}$ and $\sigma_Y^{2}(1 - \rho^{2}) = 1.458$,
whereas the exact conditional variance at $S = 7$ is $35/12 = 2.917$.  For a
Gaussian pair the conditional variance is independent of the conditioning
value; for the dice it is not, ranging from $\Var[Y \mid S = 2] = 0$ to
$\Var[Y \mid S = 7] = 35/12$.  Averaging over $S$ recovers
equation~\eqref{eq:cond-var} exactly, since $\avg{Y \mid S = s}$ is linear in
$s$.  The Gaussian expressions therefore serve below as a description of the
mechanism, and every number quoted is obtained by direct measurement rather
than from them.

\section{Application to the correlation function of a subvolume}
\label{sec:application}

\subsection{$\xi$ as a sum of squared mode amplitudes}\label{sec:xi-sum}

Expand the density contrast in the box as
\begin{equation}
  \delta(\bx) = \sum_{\bk} \delta_{\bk}\, e^{i \bk \cdot \bx},
\end{equation}
the sum running over wavevectors commensurate with the box.  Substituting twice
into the definition of $\xi$ and averaging over the box,
\begin{equation}
  \xi(r) = \frac{1}{V}\int \delta(\bx)\,\delta(\bx + \br)\, d^3x
         = \sum_{\bk}\sum_{\bk'} \delta_{\bk}\,\delta_{\bk'}\,
           e^{i \bk' \cdot \br}\,
           \frac{1}{V}\int e^{i(\bk + \bk')\cdot \bx}\, d^3x .
\end{equation}
The final integral is unity for $\bk' = -\bk$ and zero otherwise, by
orthogonality of the box modes.  Reality of $\delta$ gives $\delta_{-\bk} =
\delta_{\bk}^{*}$, so
\begin{equation}
  \xi(r) = \sum_{\bk} |\delta_{\bk}|^{2}\, e^{-i \bk \cdot \br}
         = \sum_{\bk} |\delta_{\bk}|^{2} \cos(\bk \cdot \br),
  \label{eq:xi-sum}
\end{equation}
the imaginary parts cancelling between $+\bk$ and $-\bk$.  The correlation
function is therefore a weighted sum of the squared mode amplitudes, with
weights $\cos(\bk\cdot\br)$ that suppress modes of wavelength much shorter than
$r$.  Restricting the average to a subvolume alters the weights, through the
pairs of cells the subvolume contains, but not the form of the sum.

The subvolume spans $125\Mpch$ transverse to its long axis, so the longest
transverse mode it admits has $k = 2\pi/125 \approx 0.05\ h/\mathrm{Mpc}$.
Below that wavenumber the subvolume admits only modes directed along its
length, with $k_z$ a multiple of $2\pi/1000 \approx 0.006\ h/\mathrm{Mpc}$:
approximately ten modes.  The measurement of $\xi$ near $r \sim 100\Mpch$ is
thus a sum over of order ten squared amplitudes, and matching it is the
selection of \S\ref{sec:dice} applied to those ten.  The same counting shows
why increased resolution does not help: finer sampling adds short-wavelength
modes, whereas the number of long-wavelength modes is fixed by the subvolume
size.

\subsection{The integral constraint displaces the target}\label{sec:target}

The subvolume estimator uses the mean density of the subvolume itself, as an
analysis of observational data must.  Subtracting a mean estimated from the
same data forces the volume-averaged density contrast to vanish identically,
and hence forces the integral of the measured $\xi$ over the separations the
subvolume contains to be close to zero.  An excess of pairs at small
separations is then compensated by a deficit at large separations
\citep{Peebles1980}.  The deficit is systematic: it is present in every
realization and does not diminish with the number of realizations.

Section~\ref{sec:experiment} measures it as $0.22\sigma$ at $r = 92\Mpch$.
Selecting on agreement with raw linear theory therefore requires a realization
lying $0.22\sigma$ above its own expectation in every matched bin
simultaneously, which is atypical by construction and independent of any
question of chance.

\section{Numerical experiment}\label{sec:experiment}

\subsection{Method}\label{sec:method}

We generate 600 Gaussian realizations of a $1000\Mpch$ box on a $256^3$ grid,
using the linear power spectrum from CLASS (Cosmic Linear Anisotropy Solving
System) renormalized to $\sigma_8 = 0.8$.  Only linear fields are required,
since the modes responsible for the effect grow linearly.  From each
realization we measure $\xi(r)$ in bins of $5\Mpch$ twice: over the full box,
and over the $125 \times 125 \times 1000\Mpch$ subvolume with the mean density
taken from the subvolume itself.  The ensemble mean of the full-box measurement
serves as the theoretical reference, which removes gridding and binning
differences from the comparison.

The seed search is imitated as follows.  We designate $r = 60$--$130\Mpch$,
comprising fourteen bins, as the \emph{matched range}.  For each realization we
form the sum of squared residuals between the subvolume measurement and theory
over those bins, each residual normalized by the scatter of its bin over all
600 realizations, and refer to the result as the \emph{search score}.  The
imitated search retains the $20\%$ of seeds with the lowest score, that is 120
of 600.  The procedure is implemented in
\texttt{figures/seed\_selection/plot\_seed\_conditioning.py}.

\subsection{Results}\label{sec:results}

\begin{figure}[htbp]
  \centering
  \includegraphics[width=\textwidth]{seed_conditioning.pdf}
  \caption{600 realizations of the $125 \times 125 \times 1000\Mpch$ subvolume.
    \emph{(a)} All realizations, without selection: individual measurements in
    grey, the $16$--$84$ percentile band shaded, and the ensemble mean as the
    solid curve.  The ensemble mean lies below linear theory for the reason
    given in \S\ref{sec:target}.  \emph{(b)} The realization of 600 with the
    lowest search score, plotted over the matched range and a margin of
    $10\Mpch$; it departs from theory by $0.19\sigma$ per separation.
    \emph{(c)} Scatter retained within the matched range as a function of the
    accepted fraction, together with the corresponding point for rejection on
    the mean density of the subvolume alone.  \emph{(d)} Distribution of $\xi$
    at $r = 42\Mpch$, outside the matched range, before and after selection.
    Produced by
    \texttt{figures/seed\_selection/plot\_seed\_conditioning.py}.}
  \label{fig:conditioning}
\end{figure}

\paragraph{Sample variance.}  At $r = 92\Mpch$ the subvolume measurement
scatters by $4.3\times10^{-3}$ across realizations, against $4.0\times10^{-4}$
for the full box, a ratio of $10.8$.  The volume ratio alone predicts
$\sqrt{64} = 8$; the remainder arises from the anisotropic shape of the
subvolume, which is narrow in two directions.  A survey volume of the same size
and shape is subject to the same scatter.

\paragraph{Reduction of scatter within the matched range.}
Table~\ref{tab:ratios} lists the scatter of the retained seeds divided by that
of all seeds.  Within the matched range it falls to $0.38$--$0.64$.  Any
uncertainty estimated internally from such a realization, by subvolume
splitting or an equivalent method, is correspondingly underestimated by a
factor of about two in $\sigma$, and is underestimated precisely at the
separations that were matched.

\paragraph{Propagation outside the matched range.}  The reduction extends
beyond the matched bins: the scatter is $0.67$ at $r = 58\Mpch$, $0.79$ at
$52\Mpch$ and $0.86$ at $48\Mpch$ below the range, and $0.70$ at $132\Mpch$
above it.  Ensemble means also move at separations excluded from the match, by
$-0.26\sigma$ at $r = 28\Mpch$ and $-0.24\sigma$ at $32\Mpch$, corresponding to
$2.8$ and $2.6$ times the standard error on those means.
Figure~\ref{fig:conditioning}(d) shows the distribution at $r = 42\Mpch$.

\paragraph{Displacement of the target.}  The ensemble mean of the subvolume
measurement lies $9.7\times10^{-4}$ below theory at $r = 92\Mpch$, that is
$0.22\sigma$, as anticipated in \S\ref{sec:target}.

\paragraph{The modes that are constrained.}  The combined power of the
subvolume modes with $k \leq 0.05\ h/\mathrm{Mpc}$ correlates with the search
score at $\rho = +0.38$; among the retained seeds its scatter falls to $0.67$
and its mean moves by $-0.10\sigma$.  These are the modes that set the growth
of structure within the subvolume, so the selection acts on the formation
history of the region and not only on the measured curve.

\begin{table}[htbp]
  \centering
  \caption{Effect of retaining the 120 of 600 seeds with the lowest search
    score.  Column 2 is the scatter of the retained seeds divided by that of
    all seeds; column 3 is the displacement of the ensemble mean in units of
    the scatter of all seeds; column 4 is the same displacement in units of its
    own standard error.  Values are listed at every second bin for compactness.}
  \label{tab:ratios}
  \begin{tabular}{cccc}
    \toprule
    $r$ [Mpc/$h$] & scatter retained & mean displaced & in units of its error\\
    \midrule
    \multicolumn{4}{c}{\emph{below the matched range}}\\
     28 & 1.02 & $-0.26$ & $-2.8$\\
     38 & 0.94 & $-0.19$ & $-2.1$\\
     48 & 0.86 & $-0.11$ & $-1.2$\\
     58 & 0.67 & $-0.04$ & $-0.4$\\
    \midrule
    \multicolumn{4}{c}{\emph{within the matched range}}\\
     62 & 0.56 & $+0.01$ & $+0.2$\\
     72 & 0.43 & $+0.06$ & $+0.7$\\
     82 & 0.38 & $+0.10$ & $+1.1$\\
     92 & 0.42 & $+0.13$ & $+1.4$\\
    102 & 0.41 & $+0.22$ & $+2.4$\\
    112 & 0.39 & $+0.30$ & $+3.3$\\
    122 & 0.54 & $+0.28$ & $+3.1$\\
    128 & 0.64 & $+0.23$ & $+2.5$\\
    \midrule
    \multicolumn{4}{c}{\emph{above the matched range}}\\
    132 & 0.70 & $+0.20$ & $+2.2$\\
    142 & 0.79 & $+0.09$ & $+1.0$\\
    152 & 0.83 & $+0.03$ & $+0.4$\\
    \bottomrule
  \end{tabular}
\end{table}

\section{Dependence on the tolerance}\label{sec:tolerance}

Figure~\ref{fig:conditioning}(c) varies the accepted fraction.  The scatter
retained within the matched range is

\begin{center}
\begin{tabular}{lccccccc}
  \toprule
  accepted fraction & 100\% & 80\% & 60\% & 40\% & 20\% & 10\% & 5\%\\
  scatter retained  & 1.00 & 0.76 & 0.65 & 0.55 & 0.45 & 0.37 & 0.32\\
  \bottomrule
\end{tabular}
\end{center}

\noindent
No threshold exists below which the selection is inconsequential.  Expressed as
a rejection rule rather than as a quantile: rejecting any seed whose $\xi$
departs from theory by more than $2\sigma$ at any matched separation accepts
$80.7\%$ of seeds, leaves the ensemble mean displaced by $-0.03\sigma$, and
still removes $22\%$ of the scatter.

One rule behaves differently.  Rejecting seeds whose subvolume mean density
departs from its ensemble mean by more than $2\sigma$ accepts $95.8\%$ of
seeds, displaces the mean of $\xi$ by less than $0.01\sigma$, and leaves the
scatter of $\xi$ unchanged to within the precision of 600 realizations
(measured ratio $1.01$).  The difference follows from
equation~\eqref{eq:xi-sum}: $\xi$ is a sum of squared amplitudes, whereas the
subvolume mean density is a weighted sum of the amplitudes themselves and
therefore depends on their phases.  The same set of squared amplitudes admits
both an overdense and an underdense subvolume, so constraining the mean density
within a wide interval leaves the squares nearly unconstrained.

\section{Circularity of validation by recovery}\label{sec:circular}

Consider a correlation-function estimator biased low by $10\%$, and a seed
search conducted with it.  The search will identify seeds whose reported $\xi$
lies on theory, since the distribution it samples is wide enough to contain
them, and the resulting figure will appear correct.

The relevant widths are as follows.  A $10\%$ multiplicative bias displaces the
measurement at $r = 62\Mpch$ by $0.068\sigma$, whereas the retained seeds
themselves depart from theory by $0.46\sigma$ per separation, a factor of seven
larger.  A search of this kind cannot resolve the bias.

The conclusion is a matter of inference rather than of precision.  If the seed
was chosen so that the subvolume $\xi$ agrees with theory, the observation that
the subvolume $\xi$ agrees with theory carries no information about the
estimator, the pipeline, or the simulation.  Validation of this kind requires
unselected seeds, or the full box.

A related limit applies to the procedure itself.  The closest of 600 seeds
departs from theory by $0.19\sigma$ per matched separation
(Figure~\ref{fig:conditioning}b), and the scatter within the matched range is
comparable to or larger than $\xi$ itself.  Agreement close enough to appear
exact on a figure therefore requires a search far larger than 600 seeds, at
which point \S\ref{sec:tolerance} applies with correspondingly greater force.

\section{Recommendations}\label{sec:recommendations}

\paragraph{Compare against the expectation of the estimator.}  Three
calculable effects separate the subvolume measurement from linear theory: the
window of the subvolume, the use of its own mean density (\S\ref{sec:target}),
and the grid and bin widths.  Together they place the expectation $0.22\sigma$
below theory at $r = 92\Mpch$.  Comparison against raw theory therefore
misrepresents every seed, and a search initiated from such a comparison selects
a fluctuation that compensates a systematic.  The comparison should be made
against $\avg{\hat\xi}$ for the estimator actually used.

The same consideration settles a related question.  The mean density of the
full box is available in a simulation and is not available in a survey.  If the
subvolume is used to characterize the field, normalization by the box mean
removes the effect of \S\ref{sec:target} entirely; if it represents a survey,
the subvolume mean is appropriate and the resulting suppression belongs in the
comparison curve.

\paragraph{Report the sample-variance band.}  Presenting the measurement
together with the band of Figure~\ref{fig:conditioning}(a) is sufficient
whenever the measurement lies within it.  For a subvolume representing a survey
this is also the more faithful presentation, since a survey of the same size is
subject to the same scatter, and a mock catalogue lying on theory understates
the uncertainty of a real analysis.

\paragraph{Where variance reduction is required, use a documented method.}
The objective of a realization behaving like the ensemble mean is legitimate;
the seed search is what fails.  In decreasing order of suitability for a
subvolume of this size and shape:
\begin{itemize}
  \item Fixed-amplitude initial conditions \citep{Angulo2016}, setting
        $|\delta_{\bk}| = \sqrt{P(k)}$ with random phases, optionally in
        sign-inverted pairs.  The modification is explicit and identical across
        realizations, and its effect on standard statistics has been measured
        \citep{VillaescusaNavarro2018}.  For a subvolume of this size the
        procedure removes the amplitude contribution to the scatter but not the
        contribution from the phase arrangement within the subvolume, which is
        substantial at one sixty-fourth of the volume.
  \item Constrained realizations \citep{HoffmanRibak1991}, imposing the desired
        large-scale modes directly, so that the conditioning is explicit and
        its consequences calculable.
  \item Averaging the 64 disjoint subvolumes of this shape contained in the
        box, which attains full-box precision and imposes no constraint, and
        which addresses the ensemble mean rather than an individual subvolume.
\end{itemize}

\paragraph{Where a seed must be rejected, reject on a single physical
quantity.}  If the concern is that the subvolume may be anomalously overdense
or underdense, and therefore unrepresentative for the galaxy population that
forms within it, that concern is addressed directly by rejecting on the
subvolume mean density, for which \S\ref{sec:tolerance} measures no reduction
in the scatter of $\xi$.  Rejection on the full $\xi$ curve provides the same
assurance and removes $22\%$ of the scatter.  Whichever rule is adopted, the
quantity, the tolerance, and the numbers of seeds examined and rejected should
be stated.

\paragraph{Restrict the use of a selected realization.}  Uncertainties
estimated internally at large separations, and claims that the pipeline
recovers theory (\S\ref{sec:circular}), are both compromised.  A selected
realization remains suitable for code development and timing, for figures, and
for resolution comparisons at fixed seed, in which both resolutions carry the
same selection.

\section{Conclusions}\label{sec:conclusions}

Retaining the seeds whose subvolume $\xi$ reproduces linear theory is the
selection of \S\ref{sec:dice} applied to approximately ten mode amplitudes.
Every retained realization remains a valid solution; the retained set is not a
representative sample.  For the $125 \times 125 \times 1000\Mpch$ subvolume
examined here:

\begin{enumerate}
  \item Retaining the best fifth of 600 seeds reduces the scatter of $\xi$ to
        $0.38$--$0.64$ within the matched range and to $0.67$--$0.94$ outside
        it, and displaces ensemble means at unmatched separations by up to
        $2.8$ standard errors.
  \item The reduction has no threshold.  Rejecting only the worst fifth of
        seeds removes $22\%$ of the scatter.  Rejection on the subvolume mean
        density alone removes none.
  \item The target is displaced.  The integral constraint places
        $\avg{\hat\xi}$ at $0.22\sigma$ below linear theory at $r = 92\Mpch$,
        so agreement with raw theory selects an atypical realization by
        construction.
  \item The objective is not attained.  The closest of 600 seeds departs from
        theory by $0.19\sigma$ per separation, and an estimator biased low by
        $10\%$ would be displaced by only $0.068\sigma$, so the search neither
        reproduces theory closely nor distinguishes a correct pipeline from a
        faulty one.
\end{enumerate}

The departure that motivated the proposal is in part a calculable systematic,
which belongs in the comparison curve, and in part the sample variance of a
subvolume of this size, which should be reported rather than removed.

\bibliography{references}

\end{document}
